% Generated draft: requires mathematical and source review.
\documentclass[11pt]{article}
\usepackage[margin=25mm]{geometry}
\usepackage{fontspec}
\setmainfont{Latin Modern Roman}
\usepackage{amsmath,amssymb,mathtools,unicode-math}
\setmathfont{Latin Modern Math}
\usepackage{enumitem,needspace,xurl,hyperref}
\setlength{\parindent}{0pt}
\setlength{\parskip}{5pt}
\setlength{\emergencystretch}{3em}
\newcommand{\N}{\mathbb{N}}
\newcommand{\Z}{\mathbb{Z}}
\newcommand{\Q}{\mathbb{Q}}
\newcommand{\R}{\mathbb{R}}
\newcommand{\C}{\mathbb{C}}
\newcommand{\F}{\mathbb{F}}
\newcommand{\A}{\mathbb{A}}
\newcommand{\PP}{\mathbb{P}}
\newcommand{\E}{\mathbb{E}}
\newcommand{\eps}{\varepsilon}
\newcommand{\dd}{\,\mathrm{d}}
\begin{document}
\section*{UnsolvedMath ID: 30000533}
\textbf{Source problem code:} catalogue entry
\subsection{Definitions and mathematical statement}
Let $M$ denote the arithmetic object intended by the source, and let $L_M(s)$ be its intended $L$-function. The proposed construction is to supply finite-dimensional Weil--étale cohomology spaces attached to $M$ and cup-product isomorphisms $\beta$ and $\gamma$ between the specified spaces, with a determinant comparison
\[
L_M(0)=q\,\frac{\det(\beta)}{\det(\gamma)}
\qquad\text{for some }q\in\Q^\times.
\]
Thus ``up to a rational factor'' is interpreted as up to a nonzero rational multiplicative factor. The determinant ratio requires compatible determinant-line trivializations or bases; an isomorphism between distinct vector spaces has no canonical scalar determinant without such data. A formulation modulo $\Q^\times$ additionally requires rational structures that make the ambiguity precisely rational.

The construction must use the intended Weil--étale groups and their cup products, rather than arbitrary vector spaces and arbitrarily chosen isomorphisms. The extract does not identify $M$, the class of objects quantified over, the cohomology coefficients and degrees, the cup-product classes, or the domains and codomains of $\beta,\gamma$. These missing data prevent a fully specified cohomological assertion.

The displayed expression uses the literal value $L_M(0)$. Because both maps are isomorphisms, a defined scalar determinant ratio is nonzero. Thus a zero or pole at $s=0$ would require either additional hypotheses or a leading-coefficient formulation. The source does not specify that convention, so $L_M(0)$ is not silently replaced by a leading Laurent coefficient.
\subsection{Short English statement}
Construct the intended finite-dimensional Weil--étale cohomology and cup-product isomorphisms so their determinant ratio gives $L_M(0)$ up to a nonzero rational factor. The arithmetic object, cohomology groups, determinant conventions and treatment of zeros or poles are missing from the extract.
\subsection{Sources}
\begin{itemize}
\item Supplied report: \emph{Algebraic K-Theory} (2007), Oberwolfach Reports, DOI 10.4171/owr/2006/32, pp. 1941--1990; proposed year 2006. \url{https://doi.org/10.4171/owr/2006/32}.
\item Local UnsolvedMath registry, permanent ID 30000533, catalogue entry; created 25 August 2026, updated 16 September 2026. The corrected\_verified note describes meaning-preserving editorial cleanup with no source-dependent additions. Its supplied 21 August 2026 literature triage labels this determinant comparison open and conjectural, reporting no located general proof. No additional external reference is supplied; the dated assessment has not been independently verified here.
\end{itemize}
\end{document}
